%!TEX root = ./m_main.tex

\documentclass[a4paper, 12pt]{report}

\usepackage{polyglossia}
\setdefaultlanguage{russian}

% Main font
% \setmainfont{Fira Sans}
\setmainfont{Greta Text Pro Light}
% \setmainfont{Lora}

% Monospace font
\setmonofont[Mapping=tex-text, Scale=0.85]{Fira Code}

\usepackage{amssymb}
\usepackage{mathtools}
\usepackage{amsmath}
\usepackage{amsthm}
\usepackage{eulervm}


% page props:
\usepackage[top=20mm, bottom=20mm, left=20mm, right=20mm]{geometry}
\renewcommand{\baselinestretch}{1.5}
\linespread{1.5}

% title page header
\usepackage{fancyhdr}
\usepackage{array}
\usepackage{ulem}
\newcommand{\myuline}[1]{\noindent\makebox[#1]{\hrulefill}}
\newcolumntype{L}[1]{>{\raggedright\let\newline\\\arraybackslash\hspace{0pt}}m{#1}}
\newcolumntype{C}[1]{>{\centering\let\newline\\\arraybackslash\hspace{0pt}}m{#1}}
\newcolumntype{R}[1]{>{\raggedleft\let\newline\\\arraybackslash\hspace{0pt}}m{#1}}

% \pagestyle{fancy}
% \chead{\includegraphics[width=\textwidth]{fig/aulogovec.pdf}}

% other packages

\usepackage{minibox}
\usepackage{parskip}
\usepackage{hyperref}
\usepackage{setspace}
\usepackage{multirow}
\usepackage[shortlabels]{enumitem}

% ========= Graphics packages ==========
\usepackage{graphicx}
\usepackage{import}
\usepackage{xcolor}
\usepackage{subcaption}
\usepackage{tikz}
\usepackage{tkz-berge}

% ========= Special packages ==============
\usepackage{bussproofs}
\usepackage{stmaryrd}

% ========= Listings  =========
\usepackage{minted}
\renewcommand\listingscaption{Листинг}
\usemintedstyle[python]{friendly}
\newminted{python}{
  escapeinside=~~, 
  mathescape=true, 
  fontsize=\small
}
\newminted{c}{
  escapeinside=~~, 
  mathescape=true, 
  fontsize=\small
}


% Pseudocode
\usepackage{xifthen}
\usepackage{algorithm}
\usepackage[noend]{algpseudocode}

\makeatletter
\renewcommand{\ALG@name}{Алгоритм}
\renewcommand{\listalgorithmname}{Список алгроитмов}
\newenvironment{procedure}[1]
  {\renewcommand*{\ALG@name}{Процедура}
  \algorithm\renewcommand{\thealgorithm}{\thechapter.\arabic{algorithm} #1}}
  {\endalgorithm}
\makeatother

\algrenewcommand\algorithmicrequire{\textbf{Вход:}}
\algrenewcommand\algorithmicensure{\textbf{Выход:}}
\algnewcommand\True{\textbf{true}\space}
\algnewcommand\False{\textbf{false}\space}
\algnewcommand\And{\textbf{and}\space}
\newcommand{\xfor}[3]{#1 \textbf{from} #2 \textbf{to} #3}
\newcommand{\xassign}[2]{\State #1 $\leftarrow$ #2}
\newcommand{\xstate}[1]{\State #1}
\newcommand{\xreturn}[1]{\xstate{\textbf{return} #1}}
\DeclarePairedDelimiter\ceil{\lceil}{\rceil}
\DeclarePairedDelimiter\floor{\lfloor}{\rfloor}
\newcommand{\bigO}[1]{\mathcal{O}\left(#1\right)}
\newcommand{\xqed}{\hfill $\blacksquare$}
\newcommand{\code}[1]{\colorbox{gray!15}{\footnotesize\texttt{#1}}}

% Special stuff

\usetikzlibrary{shapes,snakes}
\usetikzlibrary{arrows, petri, topaths}
\setlist[enumerate, 1]{label = (\alph*)}


\newenvironment{task}[1][]
{
  \par\medskip
  \noindent \textbf{Задание~№#1}
  \rmfamily
}
{
}

\newenvironment{solution}[1][]
{
  \par
  \noindent \textbf{\textit{Решение:~}}
  \rmfamily
}
{
  \medskip
}

% Commands (my)
\newcommand{\xparen}[1]{\left( #1 \right)}
\newcommand{\xangle}[1]{\left\langle #1 \right\rangle}
\newcommand{\xbracket}[1]{\left[ #1 \right]}
\newcommand{\xbrace}[1]{\left\lbrace #1 \right\rbrace}
\newcommand{\xret}[0]{\rightarrow}
\newcommand{\xintrp}[1]{\left\llbracket #1 \right\rrbracket}

% Math formatting
\theoremstyle{plain}
\theoremstyle{definition}
\newtheorem{lem}{Лемма}
\newtheorem{thm}{Теорема}
\newtheorem{note}{Замечание}
\newtheorem{defn}{Определение}
\newtheorem{exmp}{Пример}

\makeatletter
\def\thm@space@setup{%
  \thm@preskip=10pt
  \thm@postskip=5pt % or whatever, if you don't want them to be equal
}
\makeatother
